% !TeX program = pdflatex
\documentclass[11pt,a4paper]{article}
\usepackage[utf8]{inputenc}
\usepackage[T5]{fontenc}
\usepackage[vietnamese]{babel}
\usepackage{lmodern}
\usepackage[a4paper,left=19mm,right=19mm,top=21mm,bottom=21mm,headheight=16pt]{geometry}
\usepackage{amsmath,amssymb,array,booktabs,tabularx,longtable,ragged2e}
\usepackage[table]{xcolor}
\usepackage{tikz,listings,enumitem,fancyhdr,lastpage,needspace}
\usetikzlibrary{arrows.meta,positioning,shapes.geometric}
\usepackage[most]{tcolorbox}
\usepackage[unicode,hidelinks]{hyperref}
\usepackage{bookmark}
\definecolor{Ink}{HTML}{18384F}
\definecolor{Teal}{HTML}{087E8B}
\definecolor{Paper}{HTML}{F1F6F8}
\definecolor{Muted}{HTML}{526477}
\definecolor{MaxFill}{HTML}{D5EDF7}
\definecolor{MinFill}{HTML}{FFC2C0}
\definecolor{Green}{HTML}{276548}
\renewcommand{\tabularxcolumn}[1]{>{\RaggedRight\arraybackslash}p{#1}}
\setlength{\parindent}{0pt}\setlength{\parskip}{5pt}
\setlength{\emergencystretch}{2em}\linespread{1.02}
\setlist[enumerate]{leftmargin=6mm,itemsep=4pt,topsep=3pt}
\setlist[itemize]{leftmargin=5mm,itemsep=3pt,topsep=3pt}
\pagestyle{fancy}\fancyhf{}
\fancyhead[L]{\fontsize{8}{9.5}\selectfont\sffamily\color{Ink}BIỂU DIỄN TRI THỨC VÀ TÌM KIẾM NÂNG CAO}
\fancyhead[R]{\fontsize{8}{9.5}\selectfont\sffamily\Version}
\fancyfoot[L]{\fontsize{8.3}{10}\selectfont\color{Muted}Nguyễn Bích Vân\textbar Biểu diễn tri thức và tìm kiếm nâng cao}
\fancyfoot[R]{\footnotesize\thepage/\pageref*{LastPage}}
\renewcommand{\headrulewidth}{.4pt}
\newtcolorbox{notice}[1][]{enhanced,breakable,colback=Paper,colframe=Teal!45,boxrule=.5pt,arc=1mm,left=3mm,right=3mm,top=2mm,bottom=2mm,#1}
\newcommand{\head}[1]{\par\Needspace{3\baselineskip}\medskip{\large\sffamily\bfseries\color{Ink}#1}\par}
\newcommand{\pagehead}[1]{\clearpage\begin{tcolorbox}[colback=Ink,colframe=Ink,coltext=white,fontupper=\sffamily\large\bfseries,arc=1mm]#1\end{tcolorbox}}
\newcommand{\MaxTriangle}{\tikz[baseline=-.5ex]\node[regular polygon,regular polygon sides=3,draw=Ink,fill=MaxFill,inner sep=0pt,minimum size=3.1mm]{};}
\newcommand{\MinTriangle}{\tikz[baseline=-.5ex]\node[regular polygon,regular polygon sides=3,shape border rotate=180,draw=Ink,fill=MinFill,inner sep=0pt,minimum size=3.1mm]{};}
\lstdefinestyle{py}{language=Python,basicstyle=\ttfamily\fontsize{8.5}{10.3}\selectfont,keywordstyle=\bfseries\color{Teal},commentstyle=\color{Muted}\itshape,stringstyle=\color{Green},showstringspaces=false,keepspaces=true,columns=fullflexible,breaklines=true,frame=single,rulecolor=\color{Teal!35},backgroundcolor=\color{Paper},framesep=5pt,aboveskip=6pt,belowskip=6pt,upquote=true,tabsize=4}
\lstset{style=py}
\hypersetup{pdfauthor={Nguyễn Bích Vân},pdflang={vi-VN},pdfsubject={Biểu diễn tri thức và tìm kiếm nâng cao}}

\newcommand{\code}[1]{\nolinkurl{#1}}
\newcommand{\pts}[1]{\hfill\textbf{[#1 điểm]}}
\newcommand{\NA}{\textemdash}
\newcommand{\Version}{NHÓM 10 · MẪU BÁO CÁO}
\begin{document}\small
{\Large\sffamily\bfseries\color{Ink}01. THÔNG TIN NHÓM VÀ MÔ HÌNH}\par\medskip
BÁO CÁO GIỮA KỲ • NHÓM 10 • MINIMAX VÀ ALPHA–BETA\par\medskip
Học phần: Biểu diễn tri thức và tìm kiếm nâng cao\par Giảng viên: TS. Nguyễn Bích Vân — IAI–UET–VNU\par\medskip
Nhóm / lớp: [Điền]   Ngày giao / ngày nộp: [Điền]\par Python / hệ điều hành / máy dùng chung: [Điền]\par\medskip
Hướng dẫn: thay [Điền] bằng bài làm và kết quả thật; thêm dòng hoặc giãn ô khi cần. Mẫu trống 6 trang; báo cáo hoàn thiện 6–8 trang nội dung chính, không tính bài tính tay. Chọn Word hoặc LaTeX, không phải điền cả hai.\par\medskip
\begingroup\footnotesize\renewcommand{\arraystretch}{1.45}\begin{tabularx}{\linewidth}{XXXX}\toprule
\textbf{Thành viên / MSSV} & \textbf{Chủ trì} & \textbf{2 test phụ trách} & \textbf{Kiểm tra chéo}\\\midrule
1. [Điền] & [Điền] & [Điền] & [Điền]\\
2. [Điền] & [Điền] & [Điền] & [Điền]\\
3. [Điền] & [Điền] & [Điền] & [Điền]\\
4. [Điền] & [Điền] & [Điền] & [Điền]\\
\bottomrule\end{tabularx}\endgroup\par\medskip
\head{1. Bài toán nhóm đang giải}
\textbf{Trạng thái, hành động, chuyển trạng thái, lượt đi và kết thúc.}
Mỗi nút của cây $T_0$ là một trạng thái: gốc $R$ là trạng thái ban đầu; 14 nút trong còn lại ($A, B, A1, A2, B1, B2, A11, \dots, B22$) là trạng thái trung gian; 16 lá $L01, \dots, L16$ là trạng thái kết thúc.
Hành động tại một nút trong là chọn một trong hai con trực tiếp (trái hoặc phải).
Hàm chuyển $\mathrm{Result}(s,a)$ trả về nút con tương ứng, ví dụ $\mathrm{Result}(A, \text{phải}) = A2$.
Lượt đi luân phiên theo tầng: MAX đi ở $R$ và tầng $A1, A2, B1, B2$; MIN đi ở tầng $A, B$ và tầng $A11, \dots, B22$.
Trò chơi kết thúc khi chạm lá $L_x$ ($x \in [01, 16]$); thứ tự LR/RL chỉ đổi thứ tự xét con, không đổi tên hay vị trí lá.\par\medskip

\textbf{Utility có ý nghĩa gì? Vì sao không cộng các điểm trên đường?}
Utility là lợi ích cuối cùng của ván chơi nhìn từ phía MAX, chỉ gắn với trạng thái kết thúc (lá): MAX muốn tăng, MIN muốn giảm cùng một con số.
Các nút trong không có điểm riêng; giá trị của chúng là giá trị Minimax suy ra từ lá. Vì vậy cộng điểm dọc đường là đếm lặp và biến bài toán đối kháng thành bài toán tìm đường có tổng lớn nhất, bỏ qua việc MIN chọn ngược lại.
Utility 0 là một kết quả hợp lệ (ví dụ hòa), khác hoàn toàn với \texttt{None} (``không phải lá'').\par\medskip

\textbf{Tại sao chọn lá lớn nhất chưa phải chiến lược Minimax? Một đường chơi khác chính sách thế nào?}
Lá lớn nhất $L14 = 101$ nằm dưới $B \to B2 \to B21$. Nếu MAX chọn $B$ thì quyền quyết định tiếp theo thuộc về MIN; MIN sẽ chọn $B1$ (giá trị 4) chứ không chọn $B2$ (giá trị 100). Do đó chọn $B$ chỉ đảm bảo 4, kém hơn 10 khi chọn $A$; Minimax giả định đối thủ chơi tối ưu nên không thể ``hy vọng'' MIN đi vào $B2$.
Một đường chơi (ví dụ $R \to A \to A1 \to A11 \to L01$) chỉ là một chuỗi nước đi định sẵn, ngầm giả định đối thủ đi đúng như ta đoán.
Chính sách thì quy định nước đi của MAX tại \emph{mọi} nút MAX có thể gặp, ứng với mọi phản ứng của MIN: chọn $A$; nếu MIN chọn $A1$ thì chọn $A12$, nếu MIN chọn $A2$ thì chọn $A22$. Giá trị 10 là mức chính sách này đảm bảo dù MIN đi thế nào; đường $R \to A \to A1 \to A12 \to L04$ chỉ là đường xảy ra khi cả hai bên cùng chơi tối ưu (principal variation).\par\medskip

\textbf{Phần mã dùng nguyên, phần nhóm sửa/thêm.}
Dùng nguyên, không sửa thuật toán: \texttt{tree\_model.py} (\texttt{make\_tree}, \texttt{walk}, \texttt{validate\_tree}, \texttt{load\_cases}), \texttt{search.py} (\texttt{solve}), \texttt{data\_cases.py}, \texttt{run.py}, \texttt{tao\_du\_lieu.py}, \texttt{chay\_kiem\_thu.py}, \texttt{dong\_goi.py}.
Nhóm thêm: \texttt{test\_nhom.py} (8 test, mỗi người 2), \texttt{du\_lieu\_nhom.json} (N1, N2 sinh bằng \texttt{tao\_du\_lieu.py}), thư mục \texttt{output/} (kết quả chạy thật).\par\medskip
\pagehead{02. TRÌNH BÀY THUẬT TOÁN}
\head{2. Minimax và Alpha–Beta bằng lời của nhóm}
\textbf{Công thức truy hồi Minimax (lá / MAX / MIN).}
\[
V(u)=\begin{cases}
\mathrm{utility}(u) & \text{nếu } u \text{ là lá},\\
\max_{c \in \mathrm{con}(u)} V(c) & \text{nếu } u \text{ là nút MAX},\\
\min_{c \in \mathrm{con}(u)} V(c) & \text{nếu } u \text{ là nút MIN}.
\end{cases}
\]
Hành động tại $u$ là con $c$ đầu tiên (theo thứ tự LR hoặc RL) đạt giá trị trên; khi hòa giữ con xét trước nhờ so sánh nghiêm. Mã giả của nhóm:
\begin{lstlisting}
def minimax(u):
    if u is leaf:                      # kiem tra bang "is not None", 0 van la la
        return utility(u), None
    if u is MAX: best = -inf
    else:        best = +inf
    action = None
    for c in children(u) in order LR or RL:
        value, _ = minimax(c)
        if (u is MAX and value > best) or (u is MIN and value < best):
            best, action = value, c    # so sanh nghiem: hoa giu con truoc
    return best, action
\end{lstlisting}
\textbf{Vì sao khởi tạo $-\infty$ / $+\infty$?} Vì $\max(-\infty, x) = x$ và $\min(+\infty, x) = x$ với mọi $x$: con đầu tiên luôn ghi đè \texttt{best}, nên giá trị trả về luôn là giá trị của một con thật. Nếu khởi tạo MAX bằng 0, cây toàn điểm âm (T2, test 1) sẽ trả 0 -- con số không ứng với lá nào.\par
\textbf{Vì sao utility 0 không bị coi là ``chưa kết thúc''?} Lá được nhận biết bằng \texttt{node.utility is not None}, không phải \texttt{if node.utility}. Nếu dùng phép thử đúng/sai, lá có utility 0 bị xem như nút trong không có con và kết quả sai; \texttt{None} mới là dấu hiệu ``không phải lá'' (test 2 kiểm tra điều này).\par\medskip

Giải thích alpha, beta và fail-soft: [Người 2 điền.]\par\medskip
\begingroup\footnotesize\renewcommand{\arraystretch}{1.45}\begin{tabularx}{\linewidth}{XXX}\toprule
\textbf{Khối xử lý} & \textbf{Mã giả / giải thích của nhóm} & \textbf{Tệp và dòng}\\\midrule
Kiểm tra lá & Nếu $u$ là lá: trả \texttt{(utility, None)}. So sánh với \texttt{None} nên lá 0 vẫn là lá. & \texttt{search.py} dòng 32--36\\
Khởi tạo best/action & MAX: $\mathit{best} = -\infty$; MIN: $\mathit{best} = +\infty$; \texttt{action = None}. & \texttt{search.py} dòng 38--40\\
Gọi đệ quy & Duyệt con theo LR hoặc đảo ngược nếu RL; \texttt{score = visit(child, ...)} tương đương \texttt{minimax(c)}. & \texttt{search.py} dòng 41--43\\
Cập nhật hành động & MAX: nếu $\mathit{score} > \mathit{best}$; MIN: nếu $\mathit{score} < \mathit{best}$ thì \texttt{best, action = score, c}. So sánh nghiêm: hòa giữ con xét trước. & \texttt{search.py} dòng 44--45\\
Truyền / cập nhật cửa sổ & [Điền] & [Điền]\\
Điều kiện cắt / trả về & [Điền] & [Điền]\\
\bottomrule\end{tabularx}\endgroup\par\medskip
\head{3. Những quy tắc không được nhầm}
Tại MAX và MIN, số trả sớm là cận theo hướng nào? [Điền.]\par Vì sao so sánh chọn action nghiêm nhưng cắt dùng cả dấu bằng? [Điền.]\par\medskip
\begingroup\footnotesize\renewcommand{\arraystretch}{1.45}\begin{tabularx}{\linewidth}{XX}\toprule
\textbf{Trường} & \textbf{Nhóm hiểu và kiểm tra thế nào?}\\\midrule
visited = expanded + leaves & [Điền]\\
cutoffs / pruned\_branches & [Điền]\\
unvisited\_nodes / leaf\_order & [Điền]\\
principal\_variation / returns & [Điền]\\
\bottomrule\end{tabularx}\endgroup\par\medskip
Cận thời gian xấu nhất; bộ nhớ stack, cây sẵn có và trace: [Điền.]\par\medskip
\pagehead{03. ĐỐI CHIẾU TÍNH TAY VÀ MÃ}
PDF bài tính tay: TinhTay.pdf, trang [Điền]. Ngày ghi dự đoán trước khi chạy: [Điền].\par\medskip
\head{4. T0: Minimax và Alpha–Beta}
\begingroup\footnotesize\renewcommand{\arraystretch}{1.45}\begin{tabularx}{\linewidth}{XXX}\toprule
\textbf{Nội dung} & \textbf{Dự đoán bằng tay} & \textbf{Kết quả code / có khớp?}\\\midrule
Giá trị / hành động gốc & [Điền] & [Điền]\\
Đường chơi tối ưu từ Minimax & [Điền] & [Điền]\\
Visited / expanded / leaves & [Điền] & [Điền]\\
Gốc nhánh cắt / số cạnh & [Điền] & [Điền]\\
Lá không đọc / mọi nút không thăm & [Điền] & [Điền]\\
\bottomrule\end{tabularx}\endgroup\par\medskip
Trích một lần cắt: nút [Điền], alpha/beta [Điền], số trả [Điền], nhánh bỏ [Điền], hướng cận [Điền].\par\medskip
\head{5. T(t): thay riêng L01}
Công thức giá trị gốc theo từng khoảng t: [Điền.]\par Ngưỡng đổi đường chơi và ngưỡng cắt L04: [Điền và giải thích.]\par\medskip
\begingroup\footnotesize\renewcommand{\arraystretch}{1.45}\begin{tabularx}{\linewidth}{XXXX}\toprule
\textbf{t} & \textbf{Value / action} & \textbf{Đường chơi thật} & \textbf{Lá AB / tệp}\\\midrule
10 & [Điền] & [Điền] & [Điền]\\
11 & [Điền] & [Điền] & [Điền]\\
14 & [Điền] & [Điền] & [Điền]\\
\bottomrule\end{tabularx}\endgroup\par\medskip
\head{6. T1: kết quả điểm đúng nhưng action sai}
Giá trị thật A/B/R: [Điền]. Số B trả và hướng cận: [Điền].\par Ảnh hưởng của >= khi đổi action; phản ứng của MIN; test phát hiện: [Điền.]\par\medskip
\pagehead{04. KẾT QUẢ CHUNG VÀ DỮ LIỆU RIÊNG}
JSON/CSV chung: output/bat\_buoc.* (24 cấu hình). Lệnh và ngày chạy: [Điền].\par\medskip
\begingroup\footnotesize\renewcommand{\arraystretch}{1.45}\begin{tabularx}{\linewidth}{XXXXX}\toprule
\textbf{Cây / thứ tự} & \textbf{Value / action} & \textbf{Lá MM} & \textbf{Lá AB} & \textbf{Tiết kiệm}\\\midrule
T0/LR & [Điền] & [Điền] & [Điền] & [Điền]\\
T0/RL & [Điền] & [Điền] & [Điền] & [Điền]\\
T1/LR & [Điền] & [Điền] & [Điền] & [Điền]\\
T1/RL & [Điền] & [Điền] & [Điền] & [Điền]\\
T2/LR & [Điền] & [Điền] & [Điền] & [Điền]\\
T2/RL & [Điền] & [Điền] & [Điền] & [Điền]\\
\bottomrule\end{tabularx}\endgroup\par\medskip
Một phép tính tỷ lệ tiết kiệm: [Điền]. Một trường hợp không tiết kiệm lá: [Điền].\par\medskip
Năm mốc tham số và hai dòng affine nằm ở CSV các dòng [Điền]. Dự đoán / kết quả của u $\mapsto$ 2u-7: [Điền].\par\medskip
\head{7. Hai cây riêng N1/N2}
\begingroup\footnotesize\renewcommand{\arraystretch}{1.45}\begin{tabularx}{\linewidth}{XXX}\toprule
\textbf{Nội dung} & \textbf{N1 (16 lá, MAX)} & \textbf{N2 (8 lá, MIN)}\\\midrule
Dãy utility (hoặc dẫn JSON) & [Điền] & [Điền]\\
Seed / cách tạo & [Điền] & [Điền]\\
Giá trị / action LR & [Điền] & [Điền]\\
Giá trị / action RL & [Điền] & [Điền]\\
So sánh số lá MM và AB & [Điền] & [Điền]\\
\bottomrule\end{tabularx}\endgroup\par\medskip
Kỳ vọng tính tay độc lập của N2: [Điền / dẫn trang tính tay].\par JSON/CSV riêng: output/tu\_tao.* có 8 cấu hình; xác nhận: [Điền].\par\medskip
Đo thời gian (tùy chọn): [Không thực hiện / Điền máy, 3 lần, trung vị, trace/step đã tắt].\par\medskip
\pagehead{05. KIỂM THỬ VÀ PHÁT HIỆN SAI LẦM}
Tệp test\_nhom.py; log output/ket\_qua\_test.txt. Lệnh / ngày chạy: [Điền].\par Tổng test / đạt / lỗi thực tế: [Điền]. Không điền PASS khi chưa chạy.\par\medskip
\begingroup\footnotesize\renewcommand{\arraystretch}{1.45}\begin{tabularx}{\linewidth}{p{9mm}Xp{48mm}p{21mm}p{18mm}}\toprule
\textbf{STT} & \textbf{Test / dữ liệu} & \textbf{Kỳ vọng và cách suy ra} & \textbf{Kết quả} & \textbf{Người}\\\midrule
1 & [Điền] & [Điền] & [Điền] & [Điền]\\
2 & [Điền] & [Điền] & [Điền] & [Điền]\\
3 & [Điền] & [Điền] & [Điền] & [Điền]\\
4 & [Điền] & [Điền] & [Điền] & [Điền]\\
5 & [Điền] & [Điền] & [Điền] & [Điền]\\
6 & [Điền] & [Điền] & [Điền] & [Điền]\\
7 & [Điền] & [Điền] & [Điền] & [Điền]\\
8 & [Điền] & [Điền] & [Điền] & [Điền]\\
\bottomrule\end{tabularx}\endgroup\par\medskip
Ba test có kỳ vọng tính tay độc lập: [Điền tên test và phép tính / dẫn trang].\par Test T1 có kiểm cả value lẫn action không? [Điền].\par\medskip
\head{8. Hai lỗi nhóm đã phân tích}
\begingroup\footnotesize\renewcommand{\arraystretch}{1.45}\begin{tabularx}{\linewidth}{XXX}\toprule
\textbf{Mục} & \textbf{Lỗi 1} & \textbf{Lỗi 2}\\\midrule
Tệp / dòng / thay đổi gây lỗi & [Điền] & [Điền]\\
Dự đoán trước khi chạy & [Điền] & [Điền]\\
Dữ liệu / test phát hiện & [Điền] & [Điền]\\
Kết quả thực tế / cách sửa & [Điền] & [Điền]\\
\bottomrule\end{tabularx}\endgroup\par\medskip
Các trường hợp chưa kiểm hoặc lỗi còn tồn tại: [Điền trung thực].\par\medskip
\pagehead{06. NHẬN XÉT, TÁI TẠO VÀ BẢO VỆ}
\head{9. Ba kết luận có minh chứng}
Kết luận 1 — LR/RL và số lá: [Điền 60–90 từ, dẫn cây/cấu hình/số liệu].\par\medskip
Kết luận 2 — tham số hoặc đổi thang utility: [Điền 60–90 từ; giải thích, không chỉ chép bảng].\par\medskip
Kết luận 3 — kiểm thử và một điều hiểu rõ hơn: [Điền 60–90 từ].\par\medskip
Hai giới hạn của thực nghiệm: [Điền].\par Một cải tiến có thể khảo sát tiếp trong Minimax/Alpha–Beta: [Điền, không bắt buộc cài].\par\medskip
\head{10. Kiểm tra trước khi nộp}
\begingroup\footnotesize\renewcommand{\arraystretch}{1.45}\begin{tabularx}{\linewidth}{XX}\toprule
\textbf{Hạng mục} & \textbf{Kết quả / người kiểm}\\\midrule
Báo cáo, bài tay và 6–8 slide & [Điền]\\
24+8 cấu hình / trace T0 & [Điền]\\
8 test riêng, mỗi người 2 / log thật & [Điền]\\
README\_NOP, dữ liệu và đủ Python & [Điền]\\
ZIP giải nén chạy lại trên máy khác & [Điền]\\
Bốn người đã tập bảo vệ 15–20 phút & [Điền]\\
\bottomrule\end{tabularx}\endgroup\par\medskip
Tên ZIP: [Điền].\par Mốc ngày 5 / ngày 10 và tệp minh chứng: [Điền].\par Bốn thành viên xác nhận bài nộp và nội dung phụ trách: [Điền].\par\medskip
\end{document}
